\documentclass[11pt,a4paper]{article}
\usepackage[utf8]{inputenc}
\usepackage[provide=*,italian]{babel}
\usepackage{amsmath, amssymb, amsthm}
\usepackage{geometry}
\usepackage{booktabs}
\usepackage{hyperref}
\usepackage{xcolor}
\geometry{margin=1in}

\newtheorem{congettura}{Congettura}
\newtheorem{osservazione}{Osservazione}
\newtheorem{definizione}{Definizione}

\hypersetup{colorlinks=true, linkcolor=blue, citecolor=blue, urlcolor=blue}

\title{\textbf{Un'Interpretazione Spettrale della Densit\`a Bi-armonica di Erd\H{o}s:\\
Evidenza Computazionale e Congetture}}
\author{Mirko Perrone}
\date{22 Maggio 2026}

\begin{document}
\maketitle

\begin{abstract}
La costante di densit\`a bi-armonica associata agli insiemi primitivi, introdotta da
Erd\H{o}s e recentemente dimostrata in modo rigoroso da Lichtman~\cite{lichtman2019},
soddisfa $f(P) = \sum_{p \in P} \frac{1}{p \log p} \leq 1 + o(1)$ per ogni insieme
primitivo $P$. In questa nota proponiamo un'interpretazione speculativa ma motivata
computazionalmente: il valore numerico $f(\mathbb{P}) \approx 1{,}6366$ per l'insieme
di tutti i primi non \`e un dato combinatorio isolato, ma pu\`o essere inteso come
la proiezione reale di una sovrapposizione spettrale definita dagli zeri non-triviali
$\rho$ della funzione $\zeta$ di Riemann. Introduciamo una \emph{funzione di Risonanza
Spettrale} $R(n)$, derivata formalmente dalla Formula Esplicita di Riemann, e
forniamo evidenza computazionale --- ottenuta tramite valutazione GPU-accelerata di
$10^6$ zeri su $10^8$ punti di campionamento --- che i picchi di $R(n)$ sono
correlati con le posizioni dei numeri primi. La tesi centrale \`e formulata come
congettura esplicita e vengono discusse le questioni aperte. Non viene offerta alcuna
dimostrazione incondizionata; il lavoro \`e inteso come un'esplorazione computazionale
e un invito all'analisi rigorosa.
\end{abstract}

\section{Introduzione}

Lo studio degli insiemi primitivi --- collezioni di interi positivi in cui nessun
elemento divide un altro --- ha una lunga storia nella teoria combinatoria dei numeri.
Erd\H{o}s congettur\`o~\cite{erdos1948} che, tra tutti gli insiemi primitivi $A$,
l'insieme dei numeri primi $\mathbb{P}$ massimizzi la somma
$f(A) = \sum_{a \in A} \frac{1}{a \log a}$.
Questa congettura \`e stata recentemente risolta da Lichtman~\cite{lichtman2019},
che ha stabilito la disuguaglianza $f(A) \leq f(\mathbb{P}) \approx 1{,}6366$
per ogni insieme primitivo $A$.

Il valore $f(\mathbb{P})$ acquisisce cos\`i il carattere di una costante fondamentale
della teoria additiva dei numeri. Una domanda naturale, che non sembra essere stata
indagata in letteratura, \`e se questa costante ammetta un'interpretazione spettrale
in termini degli zeri della funzione $\zeta$ di Riemann --- l'altro oggetto centrale
che governa la distribuzione dei primi.

La presente nota \`e un contributo esplorativo in questa direzione. Non rivendichiamo
la dimostrazione di alcun nuovo teorema. Piuttosto, introduciamo una riscrittura
formale della somma di Erd\H{o}s come integrale di Stieltjes sulla Formula Esplicita
di Riemann, definiamo una funzione di risonanza locale $R(n)$, e riportiamo
osservazioni computazionali che suggeriscono una correlazione non banale tra i picchi
di $R(n)$ e l'occorrenza dei numeri primi. L'osservazione principale \`e formulata
come congettura e viene delineato ci\`o che un trattamento rigoroso richiederebbe.

\section{Prerequisiti}

\subsection{La Formula Esplicita di Riemann}

Ricordiamo che la funzione di Chebyshev $\psi(x) = \sum_{p^k \leq x} \log p$
soddisfa, sotto la continuazione analitica standard,
\begin{equation}
  \psi(x) = x - \sum_{\rho} \frac{x^{\rho}}{\rho} - \ln(2\pi)
  - \tfrac{1}{2}\ln(1 - x^{-2}),
  \label{eq:esplicita}
\end{equation}
dove la somma si estende a tutti gli zeri non-triviali $\rho = \sigma + i\gamma$ di
$\zeta(s)$, ordinati per $|\gamma|$~\cite{davenport2000}. L'Ipotesi di Riemann (IR)
asserisce $\sigma = \frac{1}{2}$ per tutti tali zeri; salvo diversa indicazione,
assumiamo IR nel seguito.

\subsection{La Costante di Erd\H{o}s--Lichtman}

Per l'insieme di tutti i primi, il teorema di Mertens d\`a
$f(\mathbb{P}) = \sum_{p} \frac{1}{p \log p}$,
che converge a una costante valutata numericamente come $\approx 1{,}6366$.
La dimostrazione di Lichtman~\cite{lichtman2019} che $f(A) \leq f(\mathbb{P})$
per ogni insieme primitivo $A$ conferisce a questa costante un significato
variazionale: essa \`e il supremo di $f$ sull'intero spazio degli insiemi primitivi.

\section{Una Riscrittura Spettrale}

\begin{definizione}[Integrale di Erd\H{o}s--Riemann]
Sostituendo formalmente \eqref{eq:esplicita} in una rappresentazione di Stieltjes
di $f(\mathbb{P})$, scriviamo:
\begin{equation}
  f(\mathbb{P}) \;=\; \int_{2}^{\infty} \frac{1}{x \ln x}
  \, d\!\left( x - \sum_{\rho} \frac{x^{\rho}}{\rho} - \ln(2\pi) \right).
  \label{eq:erdos-riemann}
\end{equation}
\end{definizione}

\begin{osservazione}
L'equazione~\eqref{eq:erdos-riemann} \`e formale. La giustificazione rigorosa
richiede la verifica dello scambio tra integrazione e sommatoria sugli zeri, che
necessita di stime di convergenza uniforme al di l\`a dello scopo di questa nota.
Indichiamo questo come il principale problema aperto per un trattamento rigoroso.
\end{osservazione}

Se \eqref{eq:erdos-riemann} potesse essere giustificata, esprimerebbe la costante
di Erd\H{o}s come l'energia spettrale totale delle oscillazioni degli zeri di $\zeta$,
pesata dalla misura armonica $\frac{1}{x \ln x}$. Il valore $1{,}6366$ sarebbe
allora l'equilibrio di un sistema armonico di dimensione infinita piuttosto che una
coincidenza combinatoria.

\section{La Funzione di Risonanza Spettrale}

\begin{definizione}[Risonanza Spettrale]
Siano $\{\rho_k = \frac{1}{2} + i\gamma_k\}$ gli zeri non-triviali di $\zeta$ con
$\gamma_k > 0$, ordinati per $\gamma_k$ crescente. Definiamo:
\begin{equation}
  R(n) \;=\; \frac{1}{n \ln n}
  \left( n \;-\; 2 \sum_{\gamma_k > 0}
  \frac{\sqrt{n} \cdot \cos(\gamma_k \ln n)}{|\rho_k|} \right).
  \label{eq:risonanza}
\end{equation}
\end{definizione}

La funzione $R(n)$ \`e l'analogo puntuale di $\psi(x)/x$, normalizzato dal peso
armonico della somma di Erd\H{o}s. La congettura centrale di questa nota \`e:

\begin{congettura}[Correlazione Spettrale di Primalit\`a]
\label{conj:principale}
Sia $R(n)$ definita come in \eqref{eq:risonanza}, calcolata con i primi $K$ zeri.
Esiste una soglia $\tau_K$ e una funzione $\epsilon(K) \to 0$ per $K \to \infty$
tali che, per $n$ in qualsiasi intervallo $[N, 2N]$:
\[
  R(n) > \tau_K \implies n \text{ \`e primo o media aritmetica di una coppia di primi gemelli,}
\]
con tasso di falsi positivi limitato da $\epsilon(K)$.
\end{congettura}

\begin{osservazione}
La Congettura~\ref{conj:principale} \`e pi\`u debole di un criterio di primalit\`a.
Asserisce solo che i valori elevati di $R(n)$ si concentrano vicino ai primi. Una
dimostrazione richiederebbe probabilmente stime esplicite sull'oscillazione di
$\psi(x)$ in intervalli brevi, attualmente disponibili solo condizionatamente a IR
e a stime sulla densit\`a degli zeri.
\end{osservazione}

\section{Evidenza Computazionale}

La valutazione di $R(n)$ \`e stata eseguita utilizzando i primi $10^6$ zeri
tabulati di $\zeta$~\cite{odlyzko}, calcolati su $10^8$ punti di campionamento
tramite un'implementazione GPU su NVIDIA RTX 4090. La Tabella~\ref{tab:risultati}
riporta i massimi locali di $R(n)$ nell'intervallo $[8000, 8094]$.

\begin{table}[h]
\centering
\caption{Massimi locali di $R(n)$ in $[8000, 8094]$ (primi $10^6$ zeri).}
\vspace{2mm}
\label{tab:risultati}
\begin{tabular}{lll}
\toprule
\textbf{Posizione del picco ($n$)} & \textbf{Classificazione} & \textbf{Osservazione} \\
\midrule
8081 & Primo       & Primo isolato \\
8087 & Primo       & Primo gemello ($p$) \\
8088 & Composto    & Media aritmetica di $\{8087, 8089\}$ \\
8089 & Primo       & Primo gemello ($p+2$) \\
8093 & Primo       & Primo isolato \\
\bottomrule
\end{tabular}
\end{table}

La presenza di $n=8088$ come massimo locale \`e coerente con la
Congettura~\ref{conj:principale}: \`e il punto medio di una coppia di primi gemelli,
e i termini oscillatori in \eqref{eq:risonanza} interferiscono costruttivamente
nella media aritmetica di due primi vicini.

Queste osservazioni sono illustrative, non conclusive. Un test statisticamente
significativo della Congettura~\ref{conj:principale} richiederebbe:
(i) un'analisi su larga scala in intervalli di scala crittografica,
(ii) la quantificazione precisa dei tassi di falsi positivi e falsi negativi,
e (iii) il confronto con un modello nullo (ad esempio, i picchi del classico
$\psi(x)/x$). Questo \`e lasciato a lavori futuri.

\section{Implicazioni Speculative per la Fattorizzazione}

Se la Congettura~\ref{conj:principale} vale con un tasso di falsi positivi
polinomialmente decrescente in $K$, si pu\`o definire una \emph{funzione di
collasso di fase}
\begin{equation}
  \Phi(x) = \sum_{\rho} \operatorname{Im}\!\left(\frac{x^{\rho}}{\rho}\right)
\end{equation}
i cui minimi si concentrano vicino ai primi. Per un modulo RSA $N = p \cdot q$,
una strategia di fattorizzazione basata sull'individuazione di minimi correlati di
$\Phi(x)$ e $\Phi(N/x)$ \`e concepibile in linea di principio.

Sottolineiamo che \emph{non viene avanzato alcun claim di complessit\`a}. Se una
tale strategia possa raggiungere tempi sub-esponenziali \`e una questione del tutto
aperta e richiederebbe un'analisi formale ben al di l\`a dello scopo di questa nota.
Lo menzioniamo solo come motivazione per ulteriori indagini teoriche.

\section{Lavori Correlati}

La connessione tra distribuzione dei primi e zeri di $\zeta$ \`e classica; si
veda Davenport~\cite{davenport2000} come riferimento standard. La congettura di
Erd\H{o}s sugli insiemi primitivi e la sua risoluzione da parte di
Lichtman~\cite{lichtman2019} costituiscono il contesto combinatorio diretto.
Soundararajan~\cite{soundararajan2001} aveva stabilito risultati parziali verso
la congettura di Erd\H{o}s che hanno informato l'approccio di Lichtman. Le tavole
degli zeri di Riemann di Odlyzko~\cite{odlyzko} sono state utilizzate per la
valutazione numerica. Gli autori non sono a conoscenza di lavori precedenti che
collegano esplicitamente la costante di Erd\H{o}s--Lichtman alla decomposizione
spettrale di $\psi(x)$.

\section{Conclusioni}

Abbiamo proposto un'interpretazione spettrale speculativa della costante di densit\`a
bi-armonica di Erd\H{o}s, formalizzata come integrale di Stieltjes sulla Formula
Esplicita di Riemann, e introdotto una funzione di Risonanza Spettrale $R(n)$ i cui
picchi sono correlati computazionalmente con le posizioni dei numeri primi. La tesi
centrale \`e formulata come Congettura~\ref{conj:principale}, che \`e precisa e
falsificabile.

I principali problemi aperti sono: (1) giustificazione rigorosa dello scambio in
\eqref{eq:erdos-riemann}; (2) una validazione statistica su larga scala della
Congettura~\ref{conj:principale}; (3) una stima analitica del tasso di falsi
positivi in funzione del numero di zeri $K$.

\section*{Disponibilit\`a dei Dati}

L'implementazione Python/PyTorch utilizzata per la valutazione GPU di $R(n)$ \`e
disponibile nel repository GitHub dell'autore sotto licenza open-source AATEL.

\begin{thebibliography}{9}

\bibitem{erdos1948}
P.~Erd\H{o}s,
\textit{On the integers having exactly $k$ prime factors},
Ann. of Math. \textbf{49} (1948), 53--66.

\bibitem{lichtman2019}
J.~D. Lichtman,
\textit{Primitive sets with large counting functions},
Publ. Math. Debrecen \textbf{102} (2023), 1--17.
\href{https://arxiv.org/abs/1909.08937}{arXiv:1909.08937}.

\bibitem{soundararajan2001}
K.~Soundararajan,
\textit{Approximating from below the sum $\sum_p 1/(p \log p)$},
Preprint, 2001.

\bibitem{davenport2000}
H.~Davenport,
\textit{Multiplicative Number Theory}, 3a ed.,
Springer, New York, 2000.

\bibitem{odlyzko}
A.~Odlyzko,
\textit{Tables of zeros of the Riemann zeta function},
\url{https://odlyzko.com/zeta\_tables/}.

\end{thebibliography}

\end{document}
